\section{BDPP's criterion}


\subsection{Movable curves}

  \begin{definition}
    A curve $C$ on a projective variety $X$ is called \emph{movable} if it belongs to an algebraic family of curves that cover $X$, i.e. there exist a variety $Y$ and a family of curves $\{C_y\}_{y \in Y}$ such that $C = C_{y_0}$ for some $y_0 \in Y$ and $\bigcup_{y \in Y} C_y = X$.
    \Yang{}
  \end{definition}



\subsection{BDPP's criterion}


    \begin{theorem}[{BDPP's criterion, ref. \cite[Theorem 0.2]{BDPP12}}]
        Let $X$ be a smooth projective variety over $\kkk$.
        Recall that $\Nef^{d-1}(X) \subset \upN^{d-1}(X)$ is defined as the dual cone of $\Psef_{d-1}(X) \subset \upN_{d-1}(X)$.
        Then we have the following equivalence of conditions describing $\Nef^{d-1}(X)$:
        \begin{enumerate}
            \item $\Nef^{d-1}(X) = \Psef^1(X)^\vee$.
            \item $\Nef^{d-1}(X) = \overline{\Mov}_1(X)$.
            \item $\Nef^{d-1}(X)$ is the closure of the cone generated by classes of form $p_*(H_1 \cdots H_{d-1})$ where $p: Y \to X$ is a birational morphism and $H_i$ are nef divisors on $Y$.
        \end{enumerate}
        \Yang{}
    \end{theorem}
